% Options for packages loaded elsewhere
\PassOptionsToPackage{unicode}{hyperref}
\PassOptionsToPackage{hyphens}{url}
%
\documentclass[
  fontsize=12pt,
  numbers=noenddot,
  parskip=half,
  listof=totoc,
  bibliography=totoc,
  headsepline=true,
  footsepline=false,
  DIV=12]{scrartcl}
\usepackage{amsmath,amssymb}
\usepackage{iftex}
\ifPDFTeX
  \usepackage[T1]{fontenc}
  \usepackage[utf8]{inputenc}
  \usepackage{textcomp} % provide euro and other symbols
\else % if luatex or xetex
  \usepackage{unicode-math} % this also loads fontspec
  \defaultfontfeatures{Scale=MatchLowercase}
  \defaultfontfeatures[\rmfamily]{Ligatures=TeX,Scale=1}
\fi
\usepackage{lmodern}
\ifPDFTeX\else
  % xetex/luatex font selection
\fi
% Use upquote if available, for straight quotes in verbatim environments
\IfFileExists{upquote.sty}{\usepackage{upquote}}{}
\IfFileExists{microtype.sty}{% use microtype if available
  \usepackage[]{microtype}
  \UseMicrotypeSet[protrusion]{basicmath} % disable protrusion for tt fonts
}{}
\makeatletter
\@ifundefined{KOMAClassName}{% if non-KOMA class
  \IfFileExists{parskip.sty}{%
    \usepackage{parskip}
  }{% else
    \setlength{\parindent}{0pt}
    \setlength{\parskip}{6pt plus 2pt minus 1pt}}
}{% if KOMA class
  \KOMAoptions{parskip=half}}
\makeatother
\usepackage{xcolor}
\usepackage{graphicx}
\makeatletter
\def\maxwidth{\ifdim\Gin@nat@width>\linewidth\linewidth\else\Gin@nat@width\fi}
\def\maxheight{\ifdim\Gin@nat@height>\textheight\textheight\else\Gin@nat@height\fi}
\makeatother
% Scale images if necessary, so that they will not overflow the page
% margins by default, and it is still possible to overwrite the defaults
% using explicit options in \includegraphics[width, height, ...]{}
\setkeys{Gin}{width=\maxwidth,height=\maxheight,keepaspectratio}
% Set default figure placement to htbp
\makeatletter
\def\fps@figure{htbp}
\makeatother
\setlength{\emergencystretch}{3em} % prevent overfull lines
\providecommand{\tightlist}{%
  \setlength{\itemsep}{0pt}\setlength{\parskip}{0pt}}
\setcounter{secnumdepth}{5}
\ifLuaTeX
\usepackage[bidi=basic]{babel}
\else
\usepackage[bidi=default]{babel}
\fi
\babelprovide[main,import]{ngerman}
% get rid of language-specific shorthands (see #6817):
\let\LanguageShortHands\languageshorthands
\def\languageshorthands#1{}
\usepackage{graphicx}
\usepackage{amsmath, amsfonts, amssymb, bbm, bm}
\usepackage[authoryear]{natbib}
\usepackage[onehalfspacing]{setspace}
\usepackage{scrlayer-scrpage}
\clearscrheadings
\automark{section}
\ihead{Fallstudien I}
\chead{}
\ohead{\headmark}
\ifoot{}
\cfoot{\pagemark}
\ofoot{}
\usepackage{booktabs}
\usepackage[hidelinks]{hyperref}
\subject{Fallstudien I}
\ifLuaTeX
  \usepackage{selnolig}  % disable illegal ligatures
\fi
\IfFileExists{bookmark.sty}{\usepackage{bookmark}}{\usepackage{hyperref}}
\IfFileExists{xurl.sty}{\usepackage{xurl}}{} % add URL line breaks if available
\urlstyle{same}
\hypersetup{
  pdftitle={Goodbye Deutschland - Eine ausgewählte Analyse des World Happiness Report},
  pdfauthor={Maximilian Petzel},
  pdflang={de},
  hidelinks,
  pdfcreator={LaTeX via pandoc}}

\title{Goodbye Deutschland - Eine ausgewählte Analyse des World
Happiness Report}
\usepackage{etoolbox}
\makeatletter
\providecommand{\subtitle}[1]{% add subtitle to \maketitle
  \apptocmd{\@title}{\par {\large #1 \par}}{}{}
}
\makeatother
\subtitle{Deskriptive Datenanalyse}
\author{Maximilian Petzel}
\date{28.10.2025}

\begin{document}
\maketitle

\thispagestyle{empty}
\newpage

\tableofcontents
\newpage

\hypertarget{einleitung}{%
\section{Einleitung}\label{einleitung}}

Der World Happiness Report ist ein Bericht, der jährlich von den
Vereinten Nationen veröffentlicht wird. Der Bericht hat das Ziel, die
Lebenszufriedenheit in verschiedenen Ländern darzustellen.

\hypertarget{erkluxe4rung-der-variablen}{%
\section{Erklärung der Variablen}\label{erkluxe4rung-der-variablen}}

\emph{(Hier deine Variablenbeschreibung einfügen.)}

\hypertarget{uxfcberschriften}{%
\section{Überschriften}\label{uxfcberschriften}}

Überschriften werden mittels Markdown-Levels definiert (z. B.
\texttt{\#}, \texttt{\#\#}, \texttt{\#\#\#}).

\hypertarget{unteruxfcberschrift}{%
\subsection{Unterüberschrift}\label{unteruxfcberschrift}}

Mit \texttt{\#\#} wird anschließend gegliedert.

\hypertarget{unterunteruxfcberschrift}{%
\subsubsection{Unterunterüberschrift}\label{unterunteruxfcberschrift}}

Noch kleinere Einheiten lassen sich mit \texttt{\#\#\#} gliedern.

\hypertarget{ohne-nummerierung}{%
\subsubsection{Ohne Nummerierung}\label{ohne-nummerierung}}

Nicht nummerierte Abschnitte kannst du als Roh-LaTeX setzen, z. B.:

\addsec{Ohne Nummerierung}

\hypertarget{textformatierung}{%
\section{Textformatierung}\label{textformatierung}}

Fließtext kann \textbf{fett}, \underline{unterstrichen}, \emph{kursiv},
\texttt{typewriter-artig} sein.

Schriftgrößen können \{\Huge Huge\}, \{\huge huge\}, \{\LARGE LARGE\},
\{\Large Large\}, \{\large large\}, \{\normalsize normalsize\},
\{\small small\}, \{\footnotesize footnotesize\},
\{\scriptsize scriptsize\} oder \{\tiny tiny\} sein.

\hypertarget{grafiken-und-tabellen}{%
\section{Grafiken und Tabellen}\label{grafiken-und-tabellen}}

Grafiken und Tabellen werden oft in LaTeX-Umgebungen gesetzt. Du kannst
in R Markdown sowohl reine Markdown-Grafiken als auch
LaTeX-\texttt{figure} verwenden.

Siehe Abbildung\textasciitilde{}\ref{grafik}, bzw.
Tabelle\textasciitilde{}\ref{tabelle}. Bei weiter entfernten Grafiken
und Tabellen, z. B. Abbildung\textasciitilde{}\ref{anhanggrafik}
(S.\textasciitilde{}\pageref{anhanggrafik}), empfiehlt sich auch die
Angabe der Seitenzahl.

\begin{figure}[h!]
  \centering
  \includegraphics[width=\textwidth]{curve}
  \caption{Jede Grafik sollte eine sinnvolle Beschreibung erhalten.}
  \label{grafik}
\end{figure}

Für die hier gezeigte Tabelle (siehe
Tabelle\textasciitilde{}\ref{tabelle}) wird das Paket \texttt{booktabs}
verwendet.

\begin{table}[h!]
  \caption{Tabellen haben eine Überschrift, die deren Inhalt kurz beschreibt.}
  \centering
  \begin{tabular}{llcccc}
    \toprule
             &         &            & \multicolumn{3}{c}{Quartile} \\\cmidrule(lr){4-6}
    Variable & Species & Mittelwert & 1. Q. & Median & 3. Q. \\
    \midrule
    Sepal Length & setosa     & 5.01 & 4.80 & 5.00 & 5.20 \\
                 & versicolor & 5.94 & 5.60 & 5.90 & 6.30 \\
                 & virginica  & 6.59 & 6.23 & 6.50 & 6.90 \\
    Sepal Width  & setosa     & 3.43 & 3.20 & 3.40 & 3.68 \\
                 & versicolor & 2.77 & 2.52 & 2.80 & 3.00 \\
                 & virginica  & 2.97 & 2.80 & 3.00 & 3.18 \\
    Petal Length & setosa     & 1.46 & 1.40 & 1.50 & 1.58 \\
                 & versicolor & 4.26 & 4.00 & 4.35 & 4.60 \\
                 & virginica  & 5.55 & 5.10 & 5.55 & 5.88 \\
    Petal Width  & setosa     & 0.25 & 0.20 & 0.20 & 0.30 \\
                 & versicolor & 1.33 & 1.20 & 1.30 & 1.50 \\
                 & virginica  & 2.03 & 1.80 & 2.00 & 2.30 \\
    \bottomrule
  \end{tabular}
  \label{tabelle}
\end{table}

\hypertarget{mathemodus}{%
\section{Mathemodus}\label{mathemodus}}

Der Mathemodus kann innerhalb eines Fließtextes gestartet und
geschlossen werden: \(X: \Omega \mapsto \mathbbm{R}\). Auch abgesetzte
Formeln können erzeugt werden: \[
  \bar{x} = \frac{1}{n} \sum_{i = 1}^n x_i .
\]

Sollen Formeln durchnummeriert werden, so kann die
\texttt{equation}-Umgebung verwendet werden: \begin{equation}
  E(X) = \int_{-\infty}^{\infty} x f(x)\, dx, \label{ewert}
\end{equation} wobei die Nummer, hier (\ref{ewert}), automatisch im Text
verwendet werden kann, wenn ein \texttt{label} vorhanden ist.

Mehrzeilige Gleichungen/Gleichungssysteme werden mit der
\texttt{align}-Umgebung ausgerichtet: \begin{align}
  1   &= 1+1 \nonumber\\
  2+2 &= 4
\end{align} Mit \texttt{nonumber} kann die Nummerierung in einzelnen
Zeilen unterdrückt werden.

Um die Durchnummerierung generell zu verhindern, kann auch
\texttt{align*} verwendet werden: \begin{align*}
  15 =&\; 5 + 4 + 3 \\
      &+ 2 + 1
\end{align*}

Sollte die Schreibweise mathematischer Formeln in LaTeX noch nicht
bekannt sein, so sei
\url{https://en.wikibooks.org/wiki/LaTeX/Mathematics} als guter
Überblick empfohlen.

\hypertarget{aufzuxe4hlungen}{%
\section{Aufzählungen}\label{aufzuxe4hlungen}}

Wichtige Umgebungen für Aufzählungen sind \texttt{enumerate},
\texttt{itemize} und \texttt{description}.

\hypertarget{section}{%
\subsection{\texorpdfstring{\texttt{enumerate}}{}}\label{section}}

\begin{enumerate}
  \item First item in a list
  \item Second item in a list \begin{enumerate}
    \item Item 1
    \item Item 2
  \end{enumerate}
  \item Third item in a list
\end{enumerate}

\hypertarget{section-1}{%
\subsection{\texorpdfstring{\texttt{itemize}}{}}\label{section-1}}

\begin{itemize}
  \item First item in a list
    \begin{itemize}
      \item First item in a list
        \begin{itemize}
          \item First item in a list
          \item Second item in a list
        \end{itemize}
      \item Second item in a list
    \end{itemize}
  \item Second item in a list
\end{itemize}

\hypertarget{section-2}{%
\subsection{\texorpdfstring{\texttt{description}}{}}\label{section-2}}

\begin{description}
  \item[GG] Grundgesamtheit
  \item[$\mathcal{N}(0,1)$] Standardnormalverteilung
  \item[u.i.v.] unabhängig identisch verteilt
  \item[ZV] Zufallsvariable
\end{description}

\hypertarget{zitate}{%
\section{Zitate}\label{zitate}}

Zitate (unter Verwendung des \texttt{natbib}-Pakets) können innerhalb
des Satzes mittels \texttt{\\citet}-Befehl eingefügt werden, wenn sie im
Literaturverzeichnis eingetragen sind, z. B. \citet{Hsu91}. Genauere
Angaben zum verwendeten Kapitel sind insbesondere bei Lehrbüchern
sinnvoll, siehe \citet[Kapitel~3, S.~222]{Hartung09}. Sollen
Literaturangaben komplett in Klammern hinter dem Satz erscheinen, so
kann \texttt{\\citep} verwendet werden \citep{Milner14}.

Jede im Literaturverzeichnis angegebene Quelle sollte im Text
referenziert werden und jede im Text angegebene Quelle muss auch im
Literaturverzeichnis vorhanden sein.

Für komplexere Möglichkeiten zur Verwaltung und Darstellung von Quellen
kann BibTeX verwendet werden.

\newpage
\appendix

\addsec{Weitere Grafiken}

\begin{figure}[h!]
  \centering
  \includegraphics[width=\textwidth]{curve}
  \caption{Die gleiche Grafik wie vorher, aber woanders.}
  \label{anhanggrafik}
\end{figure}

\begin{thebibliography}{9}
  \bibitem[Hartung et al.(2009)]{Hartung09} Hartung, J., Elpelt, B., \& Klösener, K. H. (2009). \emph{Statistik: Lehr- und Handbuch der angewandten Statistik}. Walter de Gruyter.
  \bibitem[Hsu und Wu(1991)]{Hsu91} Hsu, S. H., \& Wu, S. P. (1991). An investigation for determining the optimum length of chopsticks. \emph{Applied Ergonomics}, 22(6), 395--400.
  \bibitem[Milner und Rougier(2014)]{Milner14} Milner, K., \& Rougier, J. (2014). How to weigh a donkey in the Kenyan countryside. \emph{Significance}, 11(4), 40--43.
\end{thebibliography}

\end{document}
